\chapter{Multiple Choice Models}
\label{cap:multinomial}
% ── index ──────────────────────────────────────────────────────
\index{multinomial logit|textbf}
\index{mlogit@\texttt{mlogit()}|textbf}
\index{ordered logit|textbf}
\index{ologit@\texttt{ologit()}|textbf}
\index{ordered probit|textbf}
\index{oprobit@\texttt{oprobit()}|textbf}
\index{IIA, axiom of|textbf}
\index{independence of irrelevant alternatives|see{IIA, axiom of}}
\index{reference category}

% ─────────────────────────────────────────────────────────────────────────────
\section{Ordinal and Nominal Response}

When the outcome takes more than two categories, two structures are
possible:

\begin{itemize}
  \item \textbf{Ordinal response}: categories with a natural order (very dissatisfied,
  neutral, very satisfied; education level). Ordered models (logit/probit)
  preserve this structure.
  \item \textbf{Nominal response}: categories without order (choice of transportation
  mode; type of employment). Multinomial logit without order restriction.
\end{itemize}

\subsection*{Ordered Logit and Probit}

The model assumes a latent variable $y^* = x'\beta + \varepsilon$ and cutpoints
$c_1 < c_2 < \cdots < c_{J-1}$:

\[
  y = j \quad\Longleftrightarrow\quad c_{j-1} < y^* \leq c_j
\]

Ordered logit uses $\varepsilon \sim \mathrm{Logistic}(0,1)$; ordered probit
uses $\varepsilon \sim N(0,1)$. Both estimate $\beta$ and the cutpoints by MLE.

\subsection*{Multinomial Logit}

For $J$ alternatives, with alternative $j=J$ as base:

\[
  P(y_i = j \mid x_i) = \frac{\exp(\alpha_j + \beta_j'\,x_i)}{1 + \sum_{k=1}^{J-1} \exp(\alpha_k + \beta_k'\,x_i)}
  \quad j = 1,\ldots,J-1
\]

The coefficients $(\alpha_j, \beta_j)$ measure the effect of $x$ on the relative
odds of $j$ versus the base alternative.

% ─────────────────────────────────────────────────────────────────────────────
\section{DGP Verification}

\begin{lstlisting}[caption={DGP for ordered and multinomial models}]
seed(42)
generate df x     = rnormal()
generate df e     = (rnormal() - 0.5) * 1.7321
generate df ystar = 1.5 * x + e
// Categories 1, 2, 3 with cutoffs at 0.0 and 1.0
generate df y_ord = 1 + (ystar >= 0.0) + (ystar >= 1.0)

// Multinomial via inverse-CDF
generate df denom  = 1.0 + exp(-1.0 + 2.0*x) + exp(0.5)
generate df p1     = 1.0 / denom
generate df p12    = (1.0 + exp(-1.0 + 2.0*x)) / denom
generate df um     = rnormal()
generate df y_mlog = 1 + (um >= p1) + (um >= p12)
\end{lstlisting}

\subsection*{Ordered Logit}

\begin{lstlisting}
let m_olog = ologit(y_ord ~ x, df)
display m_olog
\end{lstlisting}

\begin{lstlisting}[language={}, basicstyle=\ttfamily\small,
  frame=none, numbers=none, backgroundcolor=\color{codbg},
  xleftmargin=0pt, xrightmargin=0pt, breaklines=false]
=========================== Ordered Logit Results ============================
N=500  Log-L=-396.4473  Pseudo R²=0.1878
-------------------------------- Coefficients --------------------------------
x        4.8137  (0.4633)  z=10.39  ***
Thresholds:  cut1=-0.1472  cut2=3.2191
\end{lstlisting}

\subsection*{Ordered Probit}

\begin{lstlisting}
let m_opro = oprobit(y_ord ~ x, df)
display m_opro
\end{lstlisting}

\begin{lstlisting}[language={}, basicstyle=\ttfamily\small,
  frame=none, numbers=none, backgroundcolor=\color{codbg},
  xleftmargin=0pt, xrightmargin=0pt, breaklines=false]
=========================== Ordered Probit Results ===========================
N=500  Log-L=-392.8369  Pseudo R²=0.1952
-------------------------------- Coefficients --------------------------------
x        2.8683  (0.2603)  z=11.02  ***
Thresholds:  cut1=-0.0768  cut2=1.9076
\end{lstlisting}

The ratio between logit and probit coefficients is $4{,}81/2{,}87 \approx 1{,}68$,
close to $\pi/\sqrt{3} \approx 1{,}81$ — the theoretical relationship between the two
scales.

\subsection*{Multinomial Logit}

\begin{lstlisting}
let m_mlog = mlogit(y_mlog ~ x, df)
display m_mlog
\end{lstlisting}

\begin{lstlisting}[language={}, basicstyle=\ttfamily\small,
  frame=none, numbers=none, backgroundcolor=\color{codbg},
  xleftmargin=0pt, xrightmargin=0pt, breaklines=false]
==================== Multinomial Logit Regression Results ====================
N=500  Log-L=-513.7171  Pseudo R²=0.0448  AIC=1035.43
------------------------------ y=1 vs base y=3 -------------------------------
const   -0.0809  (0.2024)
x       -0.9006  (0.4069)  *
------------------------------ y=2 vs base y=3 -------------------------------
const   -1.5279  (0.2542)  ***
x        2.0825  (0.4119)  ***
\end{lstlisting}

For $y=2$ versus base $y=3$: $\hat\beta_x = 2{,}08$ (true DGP: $\beta_{x,2} - \beta_{x,3} = 2{,}0$). The negative sign of $x$ in $y=1$ vs $y=3$ reflects that high $x$ increases the odds of $y=2$ (which has $e^{0{,}5}$ base probability) relative to $y=1$.

% ─────────────────────────────────────────────────────────────────────────────
\section{Syntax}

\begin{lstlisting}
ologit(y ~ x1 + x2, df)
oprobit(y ~ x1 + x2, df)
mlogit(y ~ x1 + x2, df)
mlogit(y ~ x1 + x2, df, base=1)    // explicit base alternative
\end{lstlisting}

\begin{nota}
  Multinomial logit imposes the \textbf{independence of irrelevant alternatives}
  (IIA) property: the probability ratio between two alternatives does not depend
  on the others. For panel choice data with correlated alternatives, use
  \lstinline{clogit} (ch.~\ref{cap:painel2}).
\end{nota}
